\documentclass[12pt]{article}
% Gemeinsamer Stil der Organigramme des Ministeriums fuer drakonische Sicherheit.
% Einbinden nach \documentclass mit % Gemeinsamer Stil der Organigramme des Ministeriums fuer drakonische Sicherheit.
% Einbinden nach \documentclass mit % Gemeinsamer Stil der Organigramme des Ministeriums fuer drakonische Sicherheit.
% Einbinden nach \documentclass mit \input{../ministerium_stil}.
\usepackage[a3paper,landscape,margin=12mm]{geometry}
\usepackage{fontspec}
\usepackage[ngerman]{babel}
\usepackage{etoolbox}
\usepackage{tikz}
\usetikzlibrary{calc}

\setmainfont{TeX Gyre Heros}
\pagestyle{empty}

\definecolor{nacht}{HTML}{1F2A44}
\definecolor{gold}{HTML}{B8901F}
\definecolor{iris}{HTML}{E8C547}
\definecolor{silber}{HTML}{B9C6D8}
\definecolor{papier}{HTML}{F4F2EC}
\definecolor{grau}{HTML}{5A5A5A}

\tikzset{
  hauptlinie/.style={draw=gold, line width=1.2pt},
  linie/.style={draw=nacht!60, line width=0.5pt},
  leitung/.style={draw=gold, line width=1.5pt, fill=white, inner sep=3mm,
    align=center, minimum width=70mm},
}

% #1 center of the shield in the surrounding picture's coordinates, e.g. (11,-13)
\newcommand{\wappen}[1]{%
  \begin{scope}[shift={#1}, x=10mm, y=10mm]
    \filldraw[fill=gold, draw=nacht, line width=1.2pt]
      (-1,1.1) -- (1,1.1) -- (1,0.1)
      .. controls (1,-0.7) and (0.4,-1.0) .. (0,-1.3)
      .. controls (-0.4,-1.0) and (-1,-0.7) .. (-1,0.1) -- cycle;
    \fill[nacht] (-0.75,0.05) .. controls (-0.3,0.6) and (0.3,0.6) .. (0.75,0.05)
      .. controls (0.3,-0.45) and (-0.3,-0.45) .. (-0.75,0.05) -- cycle;
    \fill[iris] (0,0.05) circle (0.29);
    \fill[nacht] (0,0.34) .. controls (0.09,0.15) and (0.09,-0.05) .. (0,-0.24)
      .. controls (-0.09,-0.05) and (-0.09,0.15) .. (0,0.34) -- cycle;
  \end{scope}%
}
.
\usepackage[a3paper,landscape,margin=12mm]{geometry}
\usepackage{fontspec}
\usepackage[ngerman]{babel}
\usepackage{etoolbox}
\usepackage{tikz}
\usetikzlibrary{calc}

\setmainfont{TeX Gyre Heros}
\pagestyle{empty}

\definecolor{nacht}{HTML}{1F2A44}
\definecolor{gold}{HTML}{B8901F}
\definecolor{iris}{HTML}{E8C547}
\definecolor{silber}{HTML}{B9C6D8}
\definecolor{papier}{HTML}{F4F2EC}
\definecolor{grau}{HTML}{5A5A5A}

\tikzset{
  hauptlinie/.style={draw=gold, line width=1.2pt},
  linie/.style={draw=nacht!60, line width=0.5pt},
  leitung/.style={draw=gold, line width=1.5pt, fill=white, inner sep=3mm,
    align=center, minimum width=70mm},
}

% #1 center of the shield in the surrounding picture's coordinates, e.g. (11,-13)
\newcommand{\wappen}[1]{%
  \begin{scope}[shift={#1}, x=10mm, y=10mm]
    \filldraw[fill=gold, draw=nacht, line width=1.2pt]
      (-1,1.1) -- (1,1.1) -- (1,0.1)
      .. controls (1,-0.7) and (0.4,-1.0) .. (0,-1.3)
      .. controls (-0.4,-1.0) and (-1,-0.7) .. (-1,0.1) -- cycle;
    \fill[nacht] (-0.75,0.05) .. controls (-0.3,0.6) and (0.3,0.6) .. (0.75,0.05)
      .. controls (0.3,-0.45) and (-0.3,-0.45) .. (-0.75,0.05) -- cycle;
    \fill[iris] (0,0.05) circle (0.29);
    \fill[nacht] (0,0.34) .. controls (0.09,0.15) and (0.09,-0.05) .. (0,-0.24)
      .. controls (-0.09,-0.05) and (-0.09,0.15) .. (0,0.34) -- cycle;
  \end{scope}%
}
.
\usepackage[a3paper,landscape,margin=12mm]{geometry}
\usepackage{fontspec}
\usepackage[ngerman]{babel}
\usepackage{etoolbox}
\usepackage{tikz}
\usetikzlibrary{calc}

\setmainfont{TeX Gyre Heros}
\pagestyle{empty}

\definecolor{nacht}{HTML}{1F2A44}
\definecolor{gold}{HTML}{B8901F}
\definecolor{iris}{HTML}{E8C547}
\definecolor{silber}{HTML}{B9C6D8}
\definecolor{papier}{HTML}{F4F2EC}
\definecolor{grau}{HTML}{5A5A5A}

\tikzset{
  hauptlinie/.style={draw=gold, line width=1.2pt},
  linie/.style={draw=nacht!60, line width=0.5pt},
  leitung/.style={draw=gold, line width=1.5pt, fill=white, inner sep=3mm,
    align=center, minimum width=70mm},
}

% #1 center of the shield in the surrounding picture's coordinates, e.g. (11,-13)
\newcommand{\wappen}[1]{%
  \begin{scope}[shift={#1}, x=10mm, y=10mm]
    \filldraw[fill=gold, draw=nacht, line width=1.2pt]
      (-1,1.1) -- (1,1.1) -- (1,0.1)
      .. controls (1,-0.7) and (0.4,-1.0) .. (0,-1.3)
      .. controls (-0.4,-1.0) and (-1,-0.7) .. (-1,0.1) -- cycle;
    \fill[nacht] (-0.75,0.05) .. controls (-0.3,0.6) and (0.3,0.6) .. (0.75,0.05)
      .. controls (0.3,-0.45) and (-0.3,-0.45) .. (-0.75,0.05) -- cycle;
    \fill[iris] (0,0.05) circle (0.29);
    \fill[nacht] (0,0.34) .. controls (0.09,0.15) and (0.09,-0.05) .. (0,-0.24)
      .. controls (-0.09,-0.05) and (-0.09,0.15) .. (0,0.34) -- cycle;
  \end{scope}%
}


% Column geometry in mm; 6 columns fill the 396 mm text width of A3 landscape.
\def\colw{61}
\def\gap{5.9}
\def\busy{-110}

\tikzset{
  stab/.style={draw=nacht!60, fill=papier, line width=0.5pt, inner sep=2mm,
    align=center, text width=112mm},
  einheit/.style={fill=nacht, text=white, inner xsep=2mm, inner ysep=3mm,
    align=center, text width=\colw mm-9mm,
    execute at begin node={\hyphenpenalty=2000 }},
  referat/.style={draw=nacht!40, fill=papier, line width=0.4pt, inner xsep=2mm,
    inner ysep=2.4mm, align=left, text width=\colw mm-13mm,
    execute at begin node={\hyphenpenalty=2000 }},
  marke/.style={draw=nacht, line width=0.6pt, minimum size=5.5mm, inner sep=0pt,
    font=\footnotesize\bfseries, text=nacht},
}

% #1 node name, #2 one of M, P, MP
\newcommand{\marken}[2]{%
  \ifstrequal{#2}{M}{\node[marke, fill=iris] at ($(#1.north east)+(-5mm,0)$) {M};}{}%
  \ifstrequal{#2}{P}{\node[marke, fill=silber] at ($(#1.north east)+(-5mm,0)$) {P};}{}%
  \ifstrequal{#2}{MP}{%
    \node[marke, fill=iris] at ($(#1.north east)+(-11.5mm,0)$) {M};%
    \node[marke, fill=silber] at ($(#1.north east)+(-5mm,0)$) {P};}{}%
}

% #1 column index
\newcommand{\spalte}[1]{%
  \pgfmathsetmacro{\spx}{#1*(\colw+\gap)}%
  \xdef\spine{\spx+1.5}%
  \coordinate (spalte-#1) at (\spx+\colw, \busy);
  \xdef\letzter{spalte-#1}%
}

% #1 id, #2 Bezeichnung, #3 Leitung (may be empty), #4 Untertitel (may be empty), #5 M, P or MP
\newcommand{\einheit}[5]{%
  \node[einheit, anchor=north east] (#1) at ($(\letzter.south east)+(0,-6mm)$)
    {\textbf{#2}%
     \ifstrempty{#3}{}{\\(#3)}%
     \ifstrempty{#4}{}{\\[2pt]{\footnotesize\color{white!75!nacht}#4}}};
  \draw[hauptlinie] (\spine,\busy) |- (#1.west);
  \marken{#1}{#5}%
  \xdef\letzter{#1}\xdef\aktiv{#1}%
}

% #1 id, #2 Bezeichnung, #3 Untertitel (may be empty)
\newcommand{\referat}[3]{%
  \node[referat, anchor=north east] (#1) at ($(\letzter.south east)+(0,-3.5mm)$)
    {#2\ifstrempty{#3}{}{\\{\footnotesize\color{grau}#3}}};
  \draw[linie] ($(\aktiv.south west)+(2mm,0)$) |- (#1.west);
  \xdef\letzter{#1}%
}

\begin{document}
\centering
\begin{tikzpicture}[x=1mm, y=1mm]
  \pgfmathsetmacro{\breite}{6*\colw+5*\gap}
  \pgfmathsetmacro{\mitte}{\breite/2}

  \wappen{(11,-15)}

  \node[anchor=north west, align=left, text=nacht] at (27,0)
    {{\large Ministerium für drakonische Sicherheit}\\[1mm]
     {\fontsize{30}{34}\selectfont\bfseries Abteilung 13}\\[2mm]
     {\large\itshape Flache Hierarchien. Kurze Entscheidungswege. Agil.}};
  \node[anchor=north east, align=right, text=grau] at (\breite,-19)
    {Stand: 22. November 2026};
  \draw[hauptlinie] (0,-33) -- (\breite,-33);

  \node[leitung, anchor=north] (leitung) at (\mitte,-40)
    {{\small\color{grau}Leitung}\\[2pt]{\Large\bfseries v.~Rossum}\\[1pt]
     {\small\color{grau}Head of Abteilung 13}};

  \node[stab, anchor=east] (buerokra) at (\mitte-15,-77)
    {{\footnotesize\color{grau}Stabsstelle}\\[1pt]Bürokratieabbau\\
     {\small Antragsstelle für Bürokratieabbau (Formular B-13a, dreifach)}};
  \node[stab, anchor=west] (enablemt) at (\mitte+15,-77)
    {{\footnotesize\color{grau}Stabsstelle}\\[1pt]Enablement\\
     {\small Empowerment, Synergien, Pizza-Freitag (Selbstzahler)}};

  \node[stab, anchor=east, fill=white, draw=gold, line width=1pt, text width=60mm] (magie) at (\mitte-15,-98)
    {magische Leitung};
  \node[stab, anchor=west, fill=white, draw=gold, line width=1pt, text width=60mm] (welt) at (\mitte+15,-98)
    {profane Leitung};
  \marken{magie}{M}
  \marken{welt}{P}

  \draw[hauptlinie] (leitung.south) -- (\mitte,\busy);
  \draw[hauptlinie] (buerokra.east) -- (\mitte,-77);
  \draw[hauptlinie] (enablemt.west) -- (\mitte,-77);
  \draw[hauptlinie] (magie.east) -- (\mitte,-98);
  \draw[hauptlinie] (welt.west) -- (\mitte,-98);
  \draw[hauptlinie] (1.5,\busy) -- (5*\colw+5*\gap+1.5,\busy);

  \spalte{0}
  \einheit{lizenz}{Magische Lizenzerteilung}{Herr Schnarch}{Bearbeitungszeit 6 bis 8 Mondzyklen}{M}
  \referat{lizenz_elementar}{Elementarmagie}{}
  \referat{lizenz_illusion}{Illusionsmagie}{}
  \referat{lizenz_beschwoer}{Beschwörung und Bannung}{}
  \referat{lizenz_weissag}{Weissagung und Hellsicht}{}
  \referat{lizenz_alchemie}{Alchemie und Tränke}{}
  \referat{lizenz_schutz}{Schutz- und Wehrzauber}{}
  \referat{lizenz_vertraut}{Vertrautenregistrierung}{}

  \spalte{1}
  \einheit{magieforensik}{Magieforensik}{v.~Wurmmoor}{Restmagie-Analyse, Fluchspuren, Ektoplasma-Labor}{M}
  \einheit{heilung}{Übernatürliche Heilung}{v.~Wurmmoor}{Bisse, Flüche, Nachsorge;\\an Vollmond geschlossen}{M}

  \spalte{2}
  \einheit{ermittlungen}{Ermittlungen}{Frau Staufer}{}{MP}
  \referat{eg1}{Ermittlungsgruppe 1}{Tagschicht}
  \referat{eg2}{Ermittlungsgruppe 2}{Tagschicht}
  \referat{eg3}{Ermittlungsgruppe 3}{Tagschicht}
  \referat{eg4}{Ermittlungsgruppe 4}{Nachtschicht}
  \referat{eg5}{Ermittlungsgruppe 5}{Bereitschaft inkl. Vollmond}

  \spalte{3}
  \einheit{beschwerden}{Beschwerden}{}{Zuständigkeitsklärung profan/magisch}{MP}
  \einheit{personal}{Personal}{}{Talent Acquisition, Einstellungstests, Onboarding für Menschen und Neumenschen}{MP}
  \einheit{controll}{Controlling}{}{Korruptionsprävention, Mittelverwendung im Sinne der Abteilung 13, KPIs und OKRs}{MP}

  \spalte{4}
  \einheit{betruung}{Betreuung}{}{Mentoring, Supervision nach Erstkontakt}{MP}
  \einheit{entwickl}{Entwicklung}{}{profan: Datenbanken, Lagebild (Thalanpir), Prozess-Mining (Seelonis), natürlich alles quelloffen und on-prem...\\[2pt]magisch: Antimagie, Portalkontrollstäbe etc.}{MP}
  \einheit{administ}{Verwaltung}{}{elektronische Aktenführung, Ausgabe elektronisch oder auf Pergament (auf Abruf)}{MP}

  \spalte{5}
  \einheit{beschaffung_profan}{Beschaffung profan}{}{Bürobedarf, Streifenwagen, Silbermunition}{P}
  \einheit{beschaffung_magisch}{Beschaffung magisch}{}{Reagenzien, Kerzen, Salz in Großgebinden, Kalteisenmunition (getrennt vom Feenpersonal lagern)}{P}

  \draw[hauptlinie] (0,-256) -- (\breite,-256);
  \node[marke, fill=iris, anchor=west] (legende-m) at (1,-263) {M};
  \node[anchor=west, text=grau] at ($(legende-m.east)+(1.5mm,0)$) {magische Leitung};
  \node[marke, fill=silber, anchor=west] (legende-p) at (45,-263) {P};
  \node[anchor=west, text=grau] at ($(legende-p.east)+(1.5mm,0)$) {profane Leitung};
  \node[anchor=east, text=nacht] at (\breite,-263)
    {\textbf{Wir stellen ein!} Offene Stellen in unserem Portal und im monatlichen Rundpergament.};
\end{tikzpicture}
\end{document}
